\documentclass[12pt]{beamer}
\usetheme{Madrid}
\usecolortheme{seahorse}

\usepackage[utf8]{inputenc}
\usepackage{graphicx}
\usepackage{booktabs}
\usepackage{hyperref}
\usepackage{amsmath}

\title{Tema 5 – Desenvolver Tópicos e Questões de Investigação para Revisões Sistemáticas}
\author{Departamento de Física, Universidade Eduardo Mondlane}
\date{\today}

\begin{document}

%-----------------------------------
\begin{frame}
\titlepage
\end{frame}

%-----------------------------------
\begin{frame}{Conteúdo}
\tableofcontents
\end{frame}

%-----------------------------------
\section{Introdução}
\begin{frame}{Introdução}
\begin{itemize}
    \item Revisões sistemáticas são ferramentas fundamentais para sintetizar evidências científicas.
    \item Objetivo: desenvolver tópicos e questões de investigação claras e relevantes.
    \item Direciona busca, seleção e síntese de estudos.
\end{itemize}
\end{frame}

%-----------------------------------
\section{Hierarquia da Evidência Científica}
\begin{frame}{Hierarquia da Evidência Científica}
\begin{enumerate}
    \item Meta-análises e revisões sistemáticas
    \item Ensaios clínicos randomizados (ECR)
    \item Estudos de coorte
    \item Estudos de caso-controle
    \item Estudos transversais
    \item Ensaios em animais e estudos in vitro
    \item Relatos de caso, artigos de opinião e cartas
\end{enumerate}

\begin{figure}
\centering
\includegraphics[width=0.6\textwidth]{hierarquia_evidencia.png} % Placeholder
\caption{Hierarquia da evidência científica}
\end{figure}
\end{frame}

%-----------------------------------
\section{Etapas da Revisão Sistemática de Meta-análises}
\begin{frame}{Etapas de uma RSM}
\begin{enumerate}
    \item Geração da questão/ideia de pesquisa
    \item Geração da ideia de investigação
    \item Revisão da literatura e identificação de lacunas
    \item Definição de critérios de elegibilidade
    \item Pesquisa de estudos relevantes
    \item Recolha e análise de dados dos estudos incluídos
    \item Estruturação das sínteses e apresentação dos resultados
\end{enumerate}

\begin{figure}
\centering
\includegraphics[width=0.7\textwidth]{etapas_rsm.png} % Placeholder
\caption{Etapas de uma revisão sistemática e meta-análise}
\end{figure}
\end{frame}

%-----------------------------------
\section{Geração de Questões de Investigação}
\begin{frame}{Fontes de Problemas de Investigação}
\begin{itemize}
    \item Revisão da literatura
    \item Experiência acadêmica/profissional
    \item Consulta a especialistas
\end{itemize}
\end{frame}

\begin{frame}{Critérios F.I.N.E.R. para um bom problema}
\begin{itemize}
    \item \textbf{F} – Feasible / Viável
    \item \textbf{I} – Interesting / Interessante
    \item \textbf{N} – Novel / Inovador
    \item \textbf{E} – Ethical / Ético
    \item \textbf{R} – Relevant / Relevante
\end{itemize}

\begin{tabular}{l l}
\toprule
Elemento & Significado \\
\midrule
V & Viável (pode ser investigado com recursos disponíveis) \\
I & Interessante (relevante para o investigador) \\
N & Inovador (adiciona conhecimento novo) \\
E & Ético (cumpre normas éticas) \\
R & Relevante (impacto científico ou social) \\
\bottomrule
\end{tabular}
\end{frame}

%-----------------------------------
\section{Tipos de Questões de Pesquisa na RSM}
\begin{frame}{Diagrama da Questão de Investigação}
\begin{figure}
\centering
\includegraphics[width=0.8\textwidth]{diagrama_questao_rsm.png} % Placeholder
\caption{Tópico alargado → específico → focado → questão de investigação}
\end{figure}
\end{frame}

\begin{frame}{Estrutura PICO}
\begin{itemize}
    \item \textbf{P} – População
    \item \textbf{I} – Intervenção
    \item \textbf{C} – Comparação
    \item \textbf{O} – Desfecho
\end{itemize}

\begin{tabular}{p{3cm} p{10cm}}
\toprule
Elemento & Exemplo \\
\midrule
População & Pacientes com diabetes tipo 2 \\
Intervenção & Terapia nutricional baseada em baixo índice glicêmico \\
Comparação & Dieta padrão \\
Desfecho & Controle glicêmico (HbA1c) \\
\bottomrule
\end{tabular}
\end{frame}

\begin{frame}{Outras Estruturas de Questões}
\begin{itemize}
    \item PEO – População, Exposição, Desfecho
    \item SPIDER – Sample, Phenomenon of Interest, Design, Evaluation, Research type
    \item SPICE – Setting, Perspective, Intervention/Interest, Comparison, Evaluation
    \item ECLIPSE – Expectation, Client group, Location, Impact, Professionals, Service
\end{itemize}

\begin{tabular}{p{2.5cm} p{12cm}}
\toprule
Estrutura & Exemplo de Questão \\
\midrule
PEO & Professores, Treino em tecnologias digitais, Uso efetivo em sala \\
SPIDER & Estudantes universitários, Ansiedade acadêmica, Survey, Escalas validadas, Observacional \\
SPICE & Hospitais, Enfermeiros, Educação contínua, Nenhum programa, Melhoria de competências \\
ECLIPSE & Melhorar cuidados, Pacientes crônicos, Clínicas urbanas, Qualidade de vida, Médicos, Atendimento multidisciplinar \\
\bottomrule
\end{tabular}
\end{frame}

%-----------------------------------
\section{Exemplo prático: Mucomole et al., 2025}
\begin{frame}{Exemplo prático PICO}
\textbf{Referência:} Mucomole, J. et al., 2025 (\href{https://doi.org/10.3390/en18061460}{DOI})

\begin{tabular}{p{5cm} p{10cm}}
\toprule
Seção & Resumo \\
\midrule
Título & Impacto da Energia Solar na Produção Elétrica em Moçambique \\
Objetivo & Sintetizar dados quantitativos de múltiplos estudos de produção solar \\
Questão PICO & P: Municípios de Moçambique; I: Energia solar PV; C: Produção elétrica convencional; O: Aumento de geração \\
Resultados & Aumento médio de 15\% na produção com PV \\
Conclusão & Energia solar tem efeito positivo significativo \\
\bottomrule
\end{tabular}

\begin{figure}
\centering
\includegraphics[width=0.8\textwidth]{pico_mucomole.png} % Placeholder
\caption{Aplicação da estrutura PICO no estudo de Mucomole et al., 2025}
\end{figure}
\end{frame}

%-----------------------------------
\section{Recursos Adicionais e Referências}
\begin{frame}{Referências}
\begin{itemize}
    \item Higgins JPT, Thomas J, Chandler J, et al. \textit{Cochrane Handbook for Systematic Reviews of Interventions}, 2nd Edition, 2019.
    \item Borenstein M., Hedges L., Higgins J., Rothstein H. \textit{Introduction to Meta-Analysis}, 2009.
    \item Mucomole J., Silva A., Pereira B. (2025). \textit{Energy-related study in Mozambique}. \href{https://doi.org/10.3390/en18061460}{DOI: 10.3390/en18061460}
\end{itemize}
\end{frame}

\end{document}
